\SetPageTheme{story}
\PageHeader
  {Structuri imbricate}
  {Cu o singura bucla numeri pe o directie. Cu doua bucle incepi sa parcurgi o suprafata intreaga.}
  {Explorare 1 din 18}

\begin{missionbox}[Unde se schimba nivelul]
Pana acum am lucrat cu stari, conditii si repetitii pe o singura axa. De aici intram in partea care duce direct spre randare: randuri, coloane, pozitii si forme.

Asta inseamna ca jocul nu mai decide doar \textit{ce} se intampla, ci si \textit{unde} se intampla.
\end{missionbox}

\BlocklyExperiment{cq-i1-grid-scan}{Parcurgere pe grila}{Leaga bucla exterioara de randuri si bucla interioara de coloane. Urmareste ordinea exacta in care sunt atinse celulele.}

\clearpage

\SetPageTheme{theory}
\PageHeader
  {Modelul de baza}
  {Exteriorul controleaza randul. Interiorul controleaza pozitia din acel rand.}
  {Explorare 2 din 18}

\begin{examplebox}[Scheletul de lucru]
\begin{lstlisting}[style=codequest]
pentru linie de la 1 la n cu pas 1 executa
    pentru coloana de la 1 la n cu pas 1 executa
        afiseaza "#"
    sfarsit
    rand nou
sfarsit
\end{lstlisting}
\end{examplebox}

\begin{conceptbox}[Cum il citim]
Nu incercam sa memoram mecanic codul. Il citim logic:
\begin{itemize}
\item bucla exterioara decide de cate ori pornim un rand nou;
\item bucla interioara decide cate pozitii completam in acel rand;
\item instructiunea \texttt{rand nou} separa randurile in output.
\end{itemize}
\end{conceptbox}

\SyntaxCheck{imbricate}

\clearpage

\SetPageTheme{guided}
\PageHeader
  {Primul patrat}
  {Inainte de diagonale si jocuri, trebuie sa vedem clar cum se naste un patrat din parcurgere.}
  {Explorare 3 din 18}

\begin{guidebox}[Trace pentru n = 4]
Deseneaza output-ul produs de modelul de baza pentru \texttt{n = 4}. Apoi raspunde:
\begin{enumerate}
\item cate randuri apar;
\item cate simboluri sunt pe fiecare rand;
\item cate celule sunt desenate in total.
\end{enumerate}
\AnswerSpace{9}
\end{guidebox}

\begin{missionbox}[De ce insistam aici]
Daca aceasta pagina este clara, tot ce urmeaza devine natural. Toate formele de mai tarziu vor fi doar variante ale aceleiasi idei: parcurgi pozitii si alegi ce afisezi in fiecare pozitie.
\end{missionbox}

\clearpage

\SetPageTheme{theory}
\PageHeader
  {Coordonate}
  {Ca sa descrii o pozitie, ai nevoie de o pereche: linie si coloana.}
  {Explorare 4 din 18}

\begin{minipage}[t]{0.48\linewidth}
\begin{checkpointbox}[Grid 5x5]
\begin{center}
\begin{tikzpicture}[scale=0.68]
  \draw[step=1cm, cqline, thin] (0,0) grid (5,5);
  \draw[cqcyan, very thick, ->] (0,0) -- (5.8,0) node[right] {\small coloana};
  \draw[cqlime, very thick, ->] (0,0) -- (0,5.8) node[above] {\small linie};
  \node[text=cqink] at (0.6,4.4) {\scriptsize (1,1)};
  \node[text=cqink] at (4.4,0.6) {\scriptsize (5,5)};
\end{tikzpicture}
\end{center}
\end{checkpointbox}
\end{minipage}
\hfill
\begin{minipage}[t]{0.48\linewidth}
\begin{missionbox}[Cum gandeste programul]
Calculatorul nu vede un patrat ca pe o singura imagine. El viziteaza pozitii una cate una. Pentru fiecare pozitie stie unde se afla si decide ce simbol sau ce culoare trebuie pusa acolo.
\end{missionbox}
\end{minipage}

\begin{semibox}[Aplicam]
Scrie trei pozitii valide dintr-o tabla 5x5 si explica in cuvinte ce reprezinta fiecare.
\AnswerSpace{5}
\end{semibox}

\clearpage

\SetPageTheme{guided}
\PageHeader
  {Diagonale}
  {Diagonalele sunt primele reguli elegante care apar atunci cand compari coordonatele.}
  {Explorare 5 din 18}

\begin{minipage}[t]{0.48\linewidth}
\begin{conceptbox}[Diagonala principala]
O pozitie este pe diagonala principala daca:
\begin{center}
\texttt{linie == coloana}
\end{center}
\end{conceptbox}
\end{minipage}
\hfill
\begin{minipage}[t]{0.48\linewidth}
\begin{conceptbox}[Diagonala secundara]
O pozitie este pe diagonala secundara daca:
\begin{center}
\texttt{linie + coloana == n + 1}
\end{center}
\end{conceptbox}
\end{minipage}

\begin{guidebox}[Aplicare]
Pentru un patrat 5x5:
\begin{enumerate}
\item scrie toate pozitiile de pe diagonala principala;
\item scrie toate pozitiile de pe diagonala secundara;
\item observa ce pozitie apartine ambelor diagonale.
\end{enumerate}
\AnswerSpace{8}
\end{guidebox}

\clearpage

\SetPageTheme{guided}
\PageHeader
  {Zone si contur}
  {Odata ce stii diagonalele, poti descrie si regiunile din jurul lor.}
  {Explorare 6 din 18}

\begin{conceptbox}[Reguli utile]
\begin{itemize}
\item deasupra diagonalei principale: \texttt{coloana > linie}
\item sub diagonala principala: \texttt{coloana < linie}
\item pe margine: \texttt{linie == 1} sau \texttt{linie == n} sau \texttt{coloana == 1} sau \texttt{coloana == n}
\end{itemize}
\end{conceptbox}

\begin{guidebox}[Doua aplicatii]
\begin{enumerate}
\item Clasifica pozitiile \texttt{(2,4), (5,1), (3,3), (1,2)} fata de diagonala principala.
\item Descrie in pseudocod cum ai afisa un patrat gol pe interior.
\end{enumerate}
\AnswerSpace{9}
\end{guidebox}

\clearpage

\SetPageTheme{guided}
\PageHeader
  {Triunghiuri si romburi}
  {Nu schimbam ideea de baza. Schimbam doar regulile dupa care limitam sau filtram pozitiile.}
  {Explorare 7 din 18}

\begin{examplebox}[Triunghi crescator]
\begin{lstlisting}[style=codequest]
pentru linie de la 1 la n cu pas 1 executa
    pentru coloana de la 1 la linie cu pas 1 executa
        afiseaza "#"
    sfarsit
    rand nou
sfarsit
\end{lstlisting}
\end{examplebox}

\begin{missionbox}[Cum gandim rombul]
Rombul se poate privi ca doua triunghiuri lipite:
\begin{enumerate}
\item in partea de sus creste numarul de simboluri;
\item in partea de jos descreste;
\item numarul de spatii se ajusteaza invers pentru a pastra forma.
\end{enumerate}
\end{missionbox}

\begin{semibox}[Aplicam]
Scrie in cuvinte ce se modifica de la un rand la altul intr-un romb: numar de spatii, numar de simboluri, punctul in care forma incepe sa se inchida.
\AnswerSpace{6}
\end{semibox}

\clearpage

\SetPageTheme{simulation}
\PageHeader
  {Simularea grilei}
  {Aici merita sa vezi efectiv traseul: ordinea vizitarii si felul in care creste numarul de celule.}
  {Explorare 8 din 18}

\SimulatorCard{cq-i1-grid-scan}{Parcurgere pe grila}{Ordinea in care sunt vizitate celulele intr-o structura imbricata si felul in care regulile de diagonala sau margine schimba marcajul final.}{Creste dimensiunea grilei si compara numarul de randuri, numarul de coloane si numarul total de celule parcurse.}

\begin{semibox}[Dupa scanare]
\begin{enumerate}
\item Ce creste liniar si ce creste mult mai repede cand maresti \texttt{n}?
\item Cum se vede in simulator diferenta dintre o regula de diagonala si o regula de margine?
\item Ce forma ai vrea sa incerci dupa ce ai inteles patratul si triunghiul?
\end{enumerate}
\AnswerSpace{6}
\end{semibox}

\clearpage

\SetPageTheme{challenge}
\PageHeader
  {Exersare pe modele}
  {Acum aplici regulile fara sa mai primesti de fiecare data drumul complet.}
  {Explorare 9 din 18}

\begin{solobox}[Set de lucru]
\begin{enumerate}
\item Scrie pseudocod pentru un patrat plin.
\item Scrie pseudocod pentru un patrat gol pe interior.
\item Scrie pseudocod pentru un patrat in care diagonala principala este marcata cu \texttt{*} si restul pozitiilor cu \texttt{.}.
\item Scrie pseudocod pentru o tabla de sah in care afisezi \texttt{X} si \texttt{O} dupa paritatea sumei coordonatelor.
\end{enumerate}
\AnswerSpace{12}
\end{solobox}

\clearpage

\SetPageTheme{story}
\PageHeader
  {De ce tocmai Pong}
  {Pong este un exemplu excelent fiindca pune la lucru aproape toate ideile invatate pana acum.}
  {Explorare 10 din 18}

\begin{missionbox}[Ce reuneste Pong]
Pong cere simultan:
\begin{itemize}
\item variabile pentru pozitia mingii si a paletelor;
\item operatori pentru miscarea lor;
\item structuri \texttt{daca} pentru margini, scor si coliziuni;
\item repetitie pentru bucla principala a jocului;
\item coordonate pentru randarea tabloului de joc.
\end{itemize}
\end{missionbox}

\begin{conceptbox}[Firul logic]
Nu construim Pong ca pe un miracol greu de atins. Il spargem in intrebari mici:
\begin{enumerate}
\item ce stare trebuie sa tin minte;
\item cum o actualizez la fiecare pas;
\item cand schimb directia sau scorul;
\item cum desenez din nou scena dupa actualizare.
\end{enumerate}
\end{conceptbox}

\clearpage

\SetPageTheme{theory}
\PageHeader
  {Starea jocului Pong}
  {Daca starea este clara, jocul incepe sa para mai putin misterios si mult mai programabil.}
  {Explorare 11 din 18}

\begin{examplebox}[Variabile esentiale]
\begin{lstlisting}[style=codequest]
mingeX <- 10
mingeY <- 6
vitezaX <- 1
vitezaY <- 1
paletaStangaY <- 5
paletaDreaptaY <- 5
scorStanga <- 0
scorDreapta <- 0
\end{lstlisting}
\end{examplebox}

\begin{missionbox}[Ce spune fiecare]
\texttt{mingeX} si \texttt{mingeY} spun unde este mingea acum. \texttt{vitezaX} si \texttt{vitezaY} spun in ce directie si cu ce pas se misca. Pozitiile paletelor spun ce poate atinge mingea. Scorurile pastreaza rezultatul vizibil pentru jucatori.
\end{missionbox}

\begin{guidebox}[Verificare]
Ce variabile s-ar schimba in urmatoarele situatii?
\begin{enumerate}
\item mingea merge spre dreapta;
\item paleta stanga urca;
\item jucatorul din dreapta castiga un punct.
\end{enumerate}
\AnswerSpace{6}
\end{guidebox}

\clearpage

\SetPageTheme{theory}
\PageHeader
  {Bucla principala}
  {Un joc precum Pong nu ruleaza o singura data. El repeta acelasi ciclu iar si iar.}
  {Explorare 12 din 18}

\begin{examplebox}[Un ciclu simplificat]
\begin{lstlisting}[style=codequest]
cat timp (jocul este activ) executa
    citeste comenzile jucatorilor
    actualizeaza pozitia mingii
    verifica margini si coliziuni
    actualizeaza scorul daca este cazul
    redeseneaza scena
sfarsit
\end{lstlisting}
\end{examplebox}

\begin{missionbox}[De ce este importanta ordinea]
Ordinea pasilor conteaza. Daca desenezi inainte sa actualizezi pozitia mingii, pe ecran ramane o informatie veche. Daca verifici scorul inainte sa vezi unde a ajuns mingea, poti rata exact momentul in care cineva castiga punctul.
\end{missionbox}

\begin{semibox}[Explica succesiunea]
Descrie in 5--6 randuri de ce acest ciclu trebuie repetat de multe ori pe secunda pentru ca jocul sa para viu.
\AnswerSpace{6}
\end{semibox}

\clearpage

\SetPageTheme{guided}
\PageHeader
  {Miscarea si coliziunile}
  {Acum vedem cum se schimba directia mingii si cand se modifica scorul.}
  {Explorare 13 din 18}

\begin{guidebox}[Trei reguli esentiale]
Completeaza in cuvinte sau pseudocod:
\begin{enumerate}
\item Ce se intampla daca mingea atinge marginea de sus sau de jos?
\item Ce se intampla daca mingea loveste o paleta?
\item Ce se intampla daca mingea trece de paleta si iese din teren?
\end{enumerate}
\AnswerSpace{10}
\end{guidebox}

\begin{conceptbox}[Observatie utila]
Multe jocuri simple par vizuale, dar in centru au doar reguli despre coordonate si stari. Asta este vestea buna: daca stii regulile, poti construi jocul pas cu pas.
\end{conceptbox}

\clearpage

\SetPageTheme{guided}
\PageHeader
  {Pong lucrat complet}
  {Punem laolalta starea, bucla, miscarea si verificarea coliziunilor.}
  {Explorare 14 din 18}

\begin{semibox}[Scenariu de urmarit]
Pornim cu:
\begin{center}
\texttt{mingeX = 10, mingeY = 6, vitezaX = 1, vitezaY = -1}
\end{center}
Urmarind trei pasi consecutivi:
\begin{enumerate}
\item scrie noua pozitie a mingii dupa fiecare pas;
\item spune daca atinge o margine;
\item spune daca ar trebui schimbata directia.
\end{enumerate}
\AnswerSpace{10}
\end{semibox}

\begin{solobox}[Mini-proiect Pong]
Descrie clar, in 6--8 randuri, cum ai organiza un mic joc Pong pe text sau pe grila: ce variabile ai folosi, ce se repeta, ce se verifica si cand se redeseneaza scena.
\AnswerSpace{8}
\end{solobox}

\clearpage

\SetPageTheme{story}
\PageHeader
  {De la Pong la TicTacToe}
  {Daca Pong iti arata miscarea continua, TicTacToe iti arata decizia pe o tabla fixa.}
  {Explorare 15 din 18}

\begin{missionbox}[De ce este urmatorul pas potrivit]
TicTacToe este o sinteza foarte buna pentru aceasta editie:
\begin{itemize}
\item ai o tabla 3x3, deci apar grile si coordonate;
\item ai mutari succesive, deci apare schimbarea starii;
\item ai verificari de castig, deci apar conditii clare;
\item ai nevoie de ordine si de reguli, dar nu ai inca fizica complicata.
\end{itemize}
\end{missionbox}

\begin{conceptbox}[Starea minima]
Trebuie sa retii:
\begin{itemize}
\item continutul celor 9 celule;
\item al cui este randul;
\item daca partida continua, este castigata sau este remiza.
\end{itemize}
\end{conceptbox}

\clearpage

\SetPageTheme{theory}
\PageHeader
  {Tabla si mutarile}
  {O mutare buna inseamna mai intai o verificare: celula este libera sau ocupata?}
  {Explorare 16 din 18}

\begin{examplebox}[Model simplificat]
\begin{lstlisting}[style=codequest]
daca (tabla[linie][coloana] este libera) atunci
    tabla[linie][coloana] <- jucatorCurent
    schimba jucatorul
altfel
    afiseaza "Mutare invalida"
sfarsit
\end{lstlisting}
\end{examplebox}

\begin{guidebox}[Aplicare]
Descrie ce se intampla in trei cazuri:
\begin{enumerate}
\item jucatorul alege o celula goala;
\item jucatorul alege o celula deja ocupata;
\item dupa mutare se completeaza o linie intreaga cu acelasi simbol.
\end{enumerate}
\AnswerSpace{8}
\end{guidebox}

\clearpage

\SetPageTheme{simulation}
\PageHeader
  {Detectia victoriei}
  {Aici se vede foarte clar legatura dintre grila, conditii si reguli de final.}
  {Explorare 17 din 18}

\SimulatorCard{cq-i1-tictactoe-board}{Tabla interactiva 3x3}{Felul in care se actualizeaza tabla dupa fiecare mutare si momentul in care o combinatie castigatoare devine adevarata.}{Simuleaza doua partide: una castigata pe linie, alta pe diagonala. Noteaza pasul exact in care conditia de victorie devine adevarata.}

\begin{semibox}[Dupa scanare]
\begin{enumerate}
\item Ce verifici prima data: liniile, coloanele sau diagonalele?
\item Cum ai detecta remiza?
\item De ce este important sa refuzi mutarile invalide inainte sa schimbi jucatorul?
\end{enumerate}
\AnswerSpace{6}
\end{semibox}

\clearpage

\SetPageTheme{challenge}
\PageHeader
  {Provocarea mare}
  {Acum ai toate piesele pentru a descrie si construi un TicTacToe clar, corect si usor de extins.}
  {Explorare 18 din 18}

\begin{bossbox}[Construieste jocul]
Construieste o versiune functionala de TicTacToe sau descrie-o complet in pseudocod, astfel incat sa aiba:
\begin{enumerate}
\item o tabla 3x3;
\item alternarea jucatorilor;
\item respingerea mutarilor invalide;
\item detectia victoriei pe linie, coloana si diagonala;
\item detectia remizei;
\item afisarea starii dupa fiecare mutare.
\end{enumerate}
\end{bossbox}

\begin{solobox}[Remix personal]
Propune o varianta alterata a jocului. Poti schimba tema, simbolurile, dimensiunea tablei sau regula de victorie. Explica in 6--8 randuri ce ai schimbat si de ce varianta ta ramane clara pentru un alt jucator.
\AnswerSpace{10}
\end{solobox}

\clearpage
